\documentclass[a4paper,12pt]{article}
\usepackage{fontspec}
\usepackage[english, russian]{babel} % Russian — последний, значит основной

\usepackage{geometry}
\geometry{left=2cm, right=2cm, top=2cm, bottom=2cm}
\usepackage{graphicx}
\usepackage{hyperref}

\setmainfont{cmunrm.otf}[
  BoldFont        = cmunbx.otf,
  ItalicFont      = cmunti.otf,
  BoldItalicFont  = cmunbi.otf
]

\usepackage{amsmath}
\usepackage{amsfonts}
\usepackage{amssymb}
\usepackage{subfiles}

\begin{document}

\section{Подготовка видеоматериала}

\subsection{Скачивание видео}

В качестве исходного материала использовался мультфильм «Вовка в тридевятом царстве» (1965).
Видео скачивалось с YouTube с помощью утилиты \texttt{yt-dlp}:

\begin{verbatim}
yt-dlp --cookies-from-browser firefox \
  --no-check-certificates \
  --js-runtimes "deno:/home/user/.deno/bin/deno" \
  --remote-components ejs:github \
  -f "bestvideo[ext=mp4]+bestaudio[ext=m4a]/best[ext=mp4]" \
  "https://youtu.be/YpygndkMSWA" \
  -o vovka.mp4
\end{verbatim}

Флаг \texttt{-{}-cookies-from-browser firefox} необходим для авторизации на YouTube.
Флаг \texttt{-{}-js-runtimes} указывает путь к JavaScript-рантайму \texttt{deno},
который требуется для обхода защиты YouTube.

\subsection{Вырезание фрагмента}

Для экспериментов был вырезан фрагмент длительностью 15 секунд:

\begin{verbatim}
ffmpeg -i vovka.mp4 -ss 00:16:10 -t 15 -an fragment.mp4
\end{verbatim}

Параметр \texttt{-ss 00:16:10} задаёт начало фрагмента,
\texttt{-t 15} — его длительность в секундах.

Параметры полученного фрагмента можно проверить командой \texttt{ffprobe}:

\begin{verbatim}
ffprobe fragment.mp4
\end{verbatim}

Результат: видео H.264, разрешение 1896×1440, частота кадров 24~fps,
аудио AAC 44100~Гц, длительность 15.01~с.



\subsection{Кодирование с предсказанием вперёд на один кадр}

Фрагмент был закодирован с использованием только P-кадров,
где каждый кадр опирается на один предыдущий (\texttt{ref=1}).
B-кадры отключены (\texttt{bframes=0}), I-кадры встречаются
только в начале (\texttt{keyint=1000}).

\begin{verbatim}
ffmpeg -i "./source_media/fragment.mp4" \
    -c:v libx265 \
    -crf 23 \
    -x265-params "keyint=1000:min-keyint=1000:bframes=0:ref=1" \
    -an \
    "./result_media/p_ref1.mp4"
\end{verbatim}

Параметр \texttt{crf=23} задаёт качество кодирования
(чем меньше значение, тем выше качество и больше размер файла).
Результирующий файл занял 10.8 МБ.



\subsection{Кодирование с предсказанием вперёд на четыре кадра}

Фрагмент был закодирован аналогично предыдущему пункту,
однако параметр \texttt{ref} увеличен до 4.

\begin{verbatim}
ffmpeg -i "./source_media/fragment.mp4" \
    -c:v libx265 \
    -crf 23 \
    -x265-params "keyint=1000:min-keyint=1000:bframes=0:ref=4" \
    -an \
    "./result_media/p_ref4.mp4"
\end{verbatim}

При \texttt{ref=4} каждый P-кадр может опираться на любой
из четырёх предыдущих кадров. При этом каждый макроблок
independently выбирает наилучший опорный кадр — один блок
может быть предсказан из кадра №1, другой из кадра №3 и т.д.
Это позволяет точнее описать движение и уменьшить остаточную
разницу, что теоретически должно дать файл меньшего размера
по сравнению с \texttt{ref=1} при том же качестве.
Результирующий файл занял 10.5 МБ.


\subsection{Кодирование с использованием P и B-кадров}

Фрагмент был закодирован с разрешёнными B-кадрами.
По сравнению с предыдущим пунктом изменён только параметр
\texttt{bframes=4} — между P-кадрами допускается до четырёх
B-кадров подряд.

\begin{verbatim}
ffmpeg -i "./source_media/fragment.mp4" \
    -c:v libx265 \
    -crf 23 \
    -x265-params "keyint=1000:min-keyint=1000:bframes=4:ref=4" \
    -an \
    "./result_media/pb.mp4"
\end{verbatim}

B-кадры используют двунаправленное предсказание — каждый блок
описывается относительно как предыдущих, так и последующих кадров.
Это даёт наиболее точное предсказание среди всех трёх вариантов,
поэтому ожидается наименьший размер файла.
Результирующий файл занял 7.4 МБ.

\subsection{Сравнение результатов}

Все результирующие файлы находятся в папке \texttt{results\_media}.
Там же находится файл \texttt{comparison.mp4}, полученный командой:

\begin{verbatim}
ffmpeg -i p_ref1.mp4 -i p_ref4.mp4 -i pb.mp4 \
    -filter_complex "\
    [0:v]drawtext=text='ref=1':fontsize=40:fontcolor=white:x=10:y=10[v0]; \
    [1:v]drawtext=text='ref=4':fontsize=40:fontcolor=white:x=10:y=10[v1]; \
    [2:v]drawtext=text='P+B':fontsize=40:fontcolor=white:x=10:y=10[v2]; \
    [v0][v1][v2]hstack=inputs=3" \
    comparison.mp4
\end{verbatim}

Файл содержит три видео, расположенных в ряд. В левом верхнем углу
каждого из них есть подпись, обозначающая тип кодирования.

Размеры полученных файлов приведены в таблице~\ref{tab:sizes}.

\begin{table}[h]
\centering
\begin{tabular}{|l|c|}
\hline
\textbf{Вариант кодирования} & \textbf{Размер файла, МБ} \\
\hline
Только P-кадры, \texttt{ref=1} & 10.8 \\
Только P-кадры, \texttt{ref=4} & 10.5 \\
P и B-кадры, \texttt{ref=4}    & 7.4  \\
\hline
\end{tabular}
\caption{Размеры файлов после кодирования}
\label{tab:sizes}
\end{table}

Результаты соответствуют теоретическим ожиданиям: добавление B-кадров
даёт наибольший выигрыш в сжатии — файл уменьшился примерно в 1.5 раза
по сравнению с вариантом \texttt{ref=1}. Разница между \texttt{ref=1}
и \texttt{ref=4} незначительна (10.8 против 10.5~МБ), что говорит о том,
что для данного материала увеличение числа опорных кадров даёт малый
эффект.

Визуальных различий между тремя вариантами обнаружено не было —
все три видео субъективно выглядят одинаково при одном значении
параметра \texttt{crf=23}.


\end{document}
